% Generated from locale/tr/content/proof-theory/sequent-calculus/quantifiers.tex; do not hand-edit.


\olfileid[tr]{pt}{seq}{qua}

\olsection{Düzenli İspatlar ve Yerine Koyma}

Niceleyici kuralları, \LeftR{\lexists} ve~\RightR{\lforall} kurallarına
uygulanan ``özdeğişken koşulları'' nedeniyle bazı güçlükler doğurur. Bu
güçlüklerin çoğu, bu çıkarımlardaki !!{constant}~$a$'nın hiçbir zaman
yeniden kullanılmaması sağlanarak giderilebilir. Şu tanımı verelim.

\begin{defn}
\Log{G1c} ya da \Log{G3c}'deki !!^a{proof}~$\pi$, aşağıdaki koşulları
sağlıyorsa \emph{düzenlidir}:
\begin{enumerate}
  \item hiçbir özdeğişken~$a$, birden fazla \LeftR{\lexists} ya da
  \RightR{\lforall} çıkarımının özdeğişkeni olarak kullanılmaz ve
  \item $a$, bir \LeftR{\lexists} ya da \RightR{\lforall} çıkarımının
  özdeğişkeniyse $\pi$ içinde yalnızca bu çıkarımın üstünde geçer.
\end{enumerate}
\end{defn}

\begin{lem}\ollabel{lem:subst-regular} $\pi$ düzenliyse ve $\Gamma \Sequent
\Delta$ ile bitiyorsa, $c$ $\pi$ içinde özdeğişken olarak kullanılmayan
bir !!a{constant} ve $t$, $\pi$'nin hiçbir özdeğişkenini içermeyen bir
terimse, $\pi$ içindeki her $c$ geçişinin $t$ ile değiştirilmesiyle elde
edilen $\Subst{\pi}{t}{c}$ düzenli bir !!{proof}{}tır. $\Subst{\pi}{t}{c}$'nin
son sekantı $\Subst{\Gamma}{t}{c} \Sequent \Subst{\Delta}{t}{c}$'dir;
burada $\Subst{\Gamma}{t}{c} = \Setabs{!A(t)}{!A(c) \in \Gamma}$.
\end{lem}

\begin{proof}
  $\pheight{\pi}$ üzerinde tümevarımla. $\pheight{\pi} = 0$ ise $\pi$
  yalnızca bir aksiyomdan oluşur. Bu durumda $\Subst{\pi}{t}{c}$ de bir
  aksiyomdur; çünkü içindeki her !!{formula}~$!A(c)$, $!A(t)$'ye dönüşür.

  $\pheight{\pi} > 0$ ise $\pi$'deki son çıkarıma göre durumları ayırırız.
  Yapısal kurallar (\Log{G1c}'de) ile önerme bağlaçlarının mantıksal
  kurallarına ilişkin durumlar kolaydır ve alıştırma olarak bırakılmıştır.
  $\pi$'nin bir niceleyici çıkarımıyla bittiği durumları ele alalım.
  \begin{enumerate}
    \item Son çıkarım $\RightR{\lforall}$ ise $\pi$ şu biçimdedir:
    \[
    \AxiomC{}
    \RightLabel{$\pi'$}
    \Deduce$\Gamma(c) \fCenter \Delta'(c), !A(a,c)$
    \RightLabel{\RightR{\lforall}}
    \UnaryInf$\Gamma(c) \fCenter \Delta'(c), \lforall[x][!A(x,c)]$
    \DisplayProof
    \]
    Tümevarım varsayımına göre $\Subst{\pi'}{t}{c}$, $\Gamma(t) \fCenter
    \Delta'(t), !A(a,t)$ sekantının düzenli bir !!{proof}{}ıdır. $a$, $\pi$'nin
    bir özdeğişkeni olduğundan $t$ içinde geçmez. Bu nedenle
    $\Subst{\pi}{t}{c}$'deki son çıkarım
    \[
    \AxiomC{}
    \RightLabel{$\Subst{\pi'}{t}{c}$}
    \Deduce$\Gamma(t) \fCenter \Delta'(t), !A(a,t)$
    \RightLabel{\RightR{\lforall}}
    \UnaryInf$\Gamma(t) \fCenter \Delta'(t), \lforall[x][!A(x,t)]$
    \DisplayProof
    \]
    doğrudur: sonuç $a$'yı içermez.
    \item Son çıkarım \Log{G3c}'de $\LeftR{\lforall}$ ise $\pi$ şu
    biçimdedir:
    \[
    \AxiomC{}
    \RightLabel{$\pi'$}
    \Deduce$!A(s(c),c), \lforall[x][!A(x,c)], \Gamma(c) \fCenter \Delta'(c)$
    \RightLabel{\LeftR{\lforall}}
    \UnaryInf$\lforall[x][!A(x,c)], \Gamma(c) \fCenter \Delta'(c)$
    \DisplayProof
    \]
    Tümevarım varsayımına göre $\Subst{\pi'}{t}{c}$, $!A(s(t),t),
    \lforall[x][!A(x,t)], \Gamma(t) \fCenter \Delta'(t)$ sekantının düzenli bir
    !!{proof}{}ıdır. $\Subst{\pi}{t}{c}$'deki son çıkarım
    \[
    \AxiomC{}
    \RightLabel{$\Subst{\pi'}{t}{c}$}
    \Deduce$!A(s(t),t), \lforall[x][!A(x,t)], \Gamma(t) \fCenter \Delta'(t)$
    \RightLabel{\LeftR{\lforall}}
    \UnaryInf$\lforall[x][!A(x,t)], \Gamma(t) \fCenter \Delta'(t)$
    \DisplayProof
    \]
    doğrudur. \Log{G1c}'deki \LeftR{\lforall} durumu benzerdir, ancak biraz
    daha basittir.
    \item Son çıkarım \RightR{\lexists} ya da \LeftR{\lexists} ise:
    alıştırma.
  \end{enumerate}
  Her durumda düzenli bir !!a{proof}{}ın (iki öncüllü kurallarda iki düzenli
  !!{proof}{}ın) sonuna bir sekant ekledik. Yeni sekant, $\pi$'nin son
  sekantından yalnızca $c$'nin $t$ terimiyle değiştirilmesi bakımından
  ayrılır. Üstelik $t$'nin $\pi$'nin hiçbir özdeğişkenini içermediği
  varsayıldığından, $\pi$ ile $\Subst{\pi}{t}{c}$ aynı özdeğişkenlere
  sahiptir ve yeni son sekant $\Subst{\pi}{t}{c}$'nin hiçbir özdeğişkenini
  içermez. Dolayısıyla $\Subst{\pi}{t}{c}$ düzenlidir.
\end{proof}

\begin{prop}
$\pi$, \Log{G1c} ya da \Log{G3c}'de bir !!a{proof} ise $\pi$ ile aynı
yapıya (özellikle aynı !!{height} değerine) sahip ve ondan yalnızca özdeğişkenlerin
değiştirilmesi bakımından ayrılan düzenli bir !!a{proof}~$\pi'$ vardır.
\end{prop}

\begin{proof}
Yalnızca bu ispat için, $\pi$'nin bir özdeğişken çıkarımına; özdeğişkeni
başka bir çıkarımın özdeğişkeni değilse ve $\pi$ içinde çıkarımın doğrudan
alt-!!{proof}{}ı dışında hiçbir yerde geçmiyorsa \emph{temiz}, aksi hâlde
\emph{kirli} diyelim. $n$, $\pi$ içindeki kirli özdeğişken çıkarımlarının
sayısı olsun. $n$ üzerinde tümevarımla $\pi$'nin istenen düzenli
!!{proof}~$\pi'$ye dönüştürülebileceğini gösteriyoruz. $n=0$ ise $\pi$
içindeki bütün özdeğişken çıkarımları temizdir; dolayısıyla $\pi$ düzenlidir
ve ispatlanacak bir şey yoktur. Aksi hâlde en üstteki bir özdeğişken
çıkarımını, örneğin \RightR{\lforall}'ı seçelim. $\pi$'nin bu çıkarımla
biten alt-!!{proof}{}ı~$\pi_1$ şu biçimdedir:
    \[
    \AxiomC{}
    \RightLabel{$\pi_2$}
    \Deduce$\Gamma \fCenter \Delta, !A(c)$
    \RightLabel{\RightR{\lforall}}
    \UnaryInf$\Gamma \fCenter \Delta, \lforall[x][!A(x)]$
    \DisplayProof
    \]
$c$ bir özdeğişken olduğundan $\Gamma$, $\Delta$ ya da
$\lforall[x][!A(x)]$ içinde geçmez. $\pi_2$ hiçbir özdeğişken çıkarımı
içermediği için düzenlidir. $a$, $\pi$ içinde geçmeyen bir !!a{constant}
olsun. \olref{lem:subst-regular} gereğince $\Subst{\pi_2}{a}{c}$,
$\Gamma \fCenter \Delta, !A(a)$'nın düzenli bir ispatıdır ve buna
$\RightR{\lforall}$ uygulayabiliriz. $\pi$ içindeki $\pi_1$'i şu
$\pi_1'$ ispatıyla değiştirirsek
    \[
    \AxiomC{}
    \RightLabel{$\Subst{\pi_2}{a}{c}$}
    \Deduce$\Gamma \fCenter \Delta, !A(a)$
    \RightLabel{\RightR{\lforall}}
    \UnaryInf$\Gamma \fCenter \Delta, \lforall[x][!A(x)]$
    \DisplayProof
    \]
kirli bir özdeğişken çıkarımını temiz bir çıkarımla değiştirmiş oluruz;
tek fark, $c$ özdeğişkeninin~$a$ ile değiştirilmesidir. Ortaya çıkan ispatta
$n-1$ kirli özdeğişken çıkarımı vardır; sonuç tümevarım varsayımından çıkar.
\end{proof}

Az önce ispatlanan önermeye açıkça başvurmayacağız; ancak ele alınan bütün
!!{proof}s{}ın düzenli olduğunu varsayacağız---düzenli değillerse
özdeğişkenlerini değiştirerek her zaman düzenli hâle getirebiliriz.

